\documentclass{article}
\usepackage[utf8]{inputenc}

\title{Summary of Rust Runtime}
\date{}

\begin{document}

\maketitle

\section{Memory Management}
\textbf{Ownership System (Compile-Time):} Rust's ownership system, including borrowing and lifetimes, is enforced entirely at compile time. This ensures memory safety and prevents data races without requiring a garbage collector or runtime checks.

\textbf{Heap Allocation (Runtime):} While the ownership system is compile-time, heap allocations (e.g., using \texttt{Box}, \texttt{Vec}, or \texttt{String}) occur at runtime. Rust’s standard library provides the allocator, but the ownership system ensures that these allocations are safe and that deallocations happen automatically when values go out of scope.

\section{Traits and Polymorphism}
\textbf{Static Dispatch (Compile-Time):} Most of Rust’s trait system is resolved at compile time. The compiler generates the appropriate code for trait methods where they are used.

\textbf{Dynamic Dispatch (Runtime):} When using trait objects (e.g., \texttt{dyn Trait}), Rust uses dynamic dispatch. This involves a vtable (virtual table) lookup at runtime to determine which method to call. This introduces some runtime overhead.

\section{Threading}
\textbf{Concurrency Primitives:} Rust’s standard library provides concurrency primitives like threads, channels, and mutexes, which rely on underlying OS features and involve runtime support.

\textbf{Thread Management:} Rust’s runtime manages thread creation, scheduling, joining, and synchronization, though the heavy lifting is done by the OS.

\section{Standard Library}
\textbf{I/O Operations:} Rust’s standard library handles I/O operations, which involve system calls and runtime support to interface with the OS.

\textbf{Error Handling:} While Rust’s \texttt{Result} and \texttt{Option} types are compile-time constructs, some runtime checks occur when handling these types.

\section{Start-up and Clean-up}
\textbf{Initialization:} Before the \texttt{main} function is called, Rust sets up a minimal runtime environment, including setting up the stack and initializing static variables.

\textbf{Clean-up:} After \texttt{main} returns, the runtime ensures any remaining resources are cleaned up, though this is minimal.

\section{Language Features}
\textbf{Closures:} Rust’s closures can capture their environment. Depending on how they capture variables, there might be a small amount of runtime overhead to manage these captures.

\section{Conclusion}
Rust's runtime is intentionally minimal, focusing on essential tasks like panic handling, dynamic dispatch, threading, I/O, and some initialization and clean-up tasks. The bulk of Rust’s safety guarantees are enforced at compile-time, leaving only a minimal runtime footprint.

\end{document}
